\documentclass[12pt,a4paper]{article}
\usepackage[utf8]{inputenc}
\usepackage{geometry}
\usepackage{listings}
\usepackage{xcolor}
\usepackage{hyperref}

\geometry{margin=1in}

% Define custom colors for code blocks
\definecolor{codegreen}{rgb}{0,0.6,0}
\definecolor{codegray}{rgb}{0.5,0.5,0.5}
\definecolor{codepurple}{rgb}{0.58,0,0.82}
\definecolor{backcolour}{rgb}{0.95,0.95,0.92}

\lstdefinestyle{mystyle}{
    backgroundcolor=\color{backcolour},   
    commentstyle=\color{codegreen},
    keywordstyle=\color{magenta},
    numberstyle=\tiny\color{codegray},
    stringstyle=\color{codepurple},
    basicstyle=\ttfamily\small,
    breakatwhitespace=false,         
    breaklines=true,                 
    captionpos=b,                    
    keepspaces=true,                 
    numbers=left,                    
    numbersep=5pt,                  
    showspaces=false,                
    showstringspaces=false,
    showtabs=false,                  
    tabsize=2
}

\lstset{style=mystyle}

\title{\textbf{My DataraFlow Cohort 2 Journey: Mastering Advanced Python (Week 2 Take-Home Exercises)}}
\author{Raji Faruq Ishola}
\date{\today}

\begin{document}

\maketitle

\noindent Hello, my wonderful network and fellow tech explorers! 🤗 

\vspace{0.3cm}
\noindent As part of my incredible learning adventure in the DataraFlow Cohort 2 program, I am thrilled to share my complete walkthrough, practical code solutions, and deep insights from our Week 2 Take-Home assignments[cite: 3]. This journey has not only sharpened my engineering skills but has also taught me how to write cleaner, more reusable, and production-ready code. My goal in sharing this comprehensive breakdown is to help anyone navigating object-oriented programming, iterators, polymorphism, and scope to gain valuable clarity and hands-on guidance! 🗺️✨

\section*{Course 1: Object-Oriented Programming (OOP) 🍪}
Object-Oriented Programming is like designing blueprints for a manufacturing factory. A \textbf{Class} serves as our master template or cookie cutter, while an \textbf{Object} represents the individual products we bake and customize from that blueprint.

\subsection*{Task 1: Constructor and String Formatting}
\textbf{Task:} Insert the missing parts to make the code return \texttt{John (36)} using a constructor and string formatting[cite: 3].

\begin{lstlisting}[language=Python]
class Person:
    # This initializes our object with unique attributes when instantiated
    def __init__(self, name, age):
        self.name = name
        self.age = age

    # This defines how our object is represented as a text string
    def __str__(self):
        return f"{self.name} ({self.age})"

# Instantiate our object with custom data
p1 = Person("John", 36)

# Print the object instance, triggering our __str__ method automatically
print(p1)
\end{lstlisting}

\noindent\textbf{Detailed Explanation and Commentary:}
\textbf{This demonstrates how we define} initialization logic using the \texttt{\_\_init\_\_} constructor and format object outputs using the special \texttt{\_\_str\_\_} method[cite: 1]. When we pass \texttt{"John"} and \texttt{36} into \texttt{Person("John", 36)}, Python binds those values to our object instance via \texttt{self}, allowing our f-string to return \texttt{John (36)} cleanly when printed[cite: 1].

\begin{center}
\fbox{\textit{[ INSERT SCREENSHOT OF COURSE 1 CODE & OUTPUT HERE ]}}
\end{center}

\subsection*{Task 2: Classes, Objects, and Properties}
\textbf{Task:} Create a class named \texttt{MyClass}, instantiate an object called \texttt{p1}, and use it to print the value of \texttt{x}[cite: 3].

\begin{lstlisting}[language=Python]
class MyClass:
    # Define a class property
    x = 5

# Instantiate our object
p1 = MyClass()

# Access and print the property using our object instance
print(p1.x)
\end{lstlisting}

\noindent\textbf{Detailed Explanation and Commentary:}
\textbf{This illustrates how an object} accesses variables stored inside its parent class blueprint[cite: 1]. By writing \texttt{p1 = MyClass()}, we bring the template to life, and \texttt{p1.x} reaches inside to fetch the stored value \texttt{5}[cite: 1].

\section*{Course 2: Inheritance 🧬✨}
Inheritance allows us to pass down traits from a parent class to a child class, keeping our codebase DRY (Don't Repeat Yourself)[cite: 2]. Think of it like a family recipe book where children inherit master recipes and add their own signature twists.

\subsection*{Task 1: Animal Subclass Sound Overriding}
\textbf{Task:} Create a base class \texttt{Animal} with attributes \texttt{name} and \texttt{sound}, and a method \texttt{make_sound()}. Create subclasses \texttt{Dog} and \texttt{Cat} that override \texttt{make_sound()} to print \texttt{"Woof!"} and \texttt{"Meow!"}[cite: 3].

\begin{lstlisting}[language=Python]
class Animal:
    # This sets up our base animal attributes
    def __init__(self, name, sound):
        self.name = name
        self.sound = sound

    def make_sound(self):
        print(self.sound)

class Dog(Animal):
    def __init__(self, name):
        # This uses super() to inherit initialization from the parent Animal class
        super().__init__(name, "Woof!")

# Instantiate and test our subclass
dog = Dog("Bingo")
dog.make_sound() # Output: Woof!
\end{lstlisting}

\noindent\textbf{Detailed Explanation and Commentary:}
\textbf{This illustrates how \texttt{super().\_\_init\_\_()}} bridges our child class to the parent class, letting us inherit base setup logic without rewriting code[cite: 2]. By extending \texttt{Animal}, our \texttt{Dog} class gains all parent capabilities while tailoring its own unique sound behavior[cite: 2].

\subsection*{Task 2: Vehicle Hierarchy and Method Overriding}
\textbf{Task:} Create a base class \texttt{Vehicle} with attributes \texttt{brand} and \texttt{model}, and a method \texttt{display_info()}. Create a subclass \texttt{Car} that adds \texttt{doors} and overrides \texttt{display_info()}[cite: 3].

\begin{lstlisting}[language=Python]
class Vehicle:
    def __init__(self, brand, model):
        self.brand = brand
        self.model = model

    def display_info(self):
        print(f"Vehicle: {self.brand} {self.model}")

class Car(Vehicle):
    def __init__(self, brand, model, doors):
        super().__init__(brand, model)
        self.doors = doors

    def display_info(self):
        print(f"Car: {self.brand} {self.model}, Doors: {self.doors}")

# --- Execution ---
my_car = Car("Toyota", "Camry", 4)
my_car.display_info()
\end{lstlisting}

\noindent\textbf{Detailed Explanation and Commentary:}
\textbf{This demonstrates method overriding} in inheritance, where a child class modifies an inherited method to present specialized information unique to its type[cite: 2].

\section*{Course 3: Iterators 🌀✨}
An iterator acts like a bookmark in a massive book, allowing us to traverse through elements sequentially one item at a time without loading an entire collection into computer memory all at once.

\subsection*{Task 1: Custom CountDown Iterator}
\textbf{Task:} Create a custom iterator class \texttt{CountDown} that takes a number \texttt{n} and counts down to 1 using \texttt{\_\_iter\_\_()} and \texttt{\_\_next\_\_()}[cite: 3].

\begin{lstlisting}[language=Python]
class CountDown:
    def __init__(self, n):
        self.n = n

    def __iter__(self):
        # This returns the iterator object itself to initiate the protocol
        return self

    def __next__(self):
        # This yields the next sequential value or stops the iteration
        if self.n > 0:
            val = self.n
            self.n -= 1
            return val
        else:
            raise StopIteration

# --- Execution ---
for num in CountDown(3):
    print(num)
\end{lstlisting}

\noindent\textbf{Detailed Explanation and Commentary:}
\textbf{This demonstrates how custom iterators} maintain internal state across loop cycles[cite: 3]. Every time the \texttt{for} loop requests an item, \texttt{\_\_next\_\_()} executes, decrementing our counter until it hits zero and raises \texttt{StopIteration} to exit cleanly[cite: 3].

\subsection*{Task 2: List Iterator with Exception Handling}
\textbf{Task:} Take a list of numbers \texttt{[10, 20, 30, 40, 50]}, convert it into an iterator using \texttt{iter()}, and retrieve each element using \texttt{next()} while handling \texttt{StopIteration} gracefully[cite: 3].

\begin{lstlisting}[language=Python]
numbers = [10, 20, 30, 40, 50]
my_iter = iter(numbers)

while True:
    try:
        # Retrieve items sequentially one by one
        print(next(my_iter))
    except StopIteration:
        print("Reached the end of the iterator safely!")
        break
\end{lstlisting}

\noindent\textbf{Detailed Explanation and Commentary:}
\textbf{This illustrates how we process} sequences element-by-element using explicit \texttt{next()} calls and catch the \texttt{StopIteration} exception to terminate our loop without crashing[cite: 3].

\section*{Course 4: Polymorphism 🎭✨}
Polymorphism means "many forms". It allows different classes to share identical method names while executing completely customized behaviors behind the scenes.

\subsection*{Task 1: Bird Language and Consumer Function}
\textbf{Task:} Create a base class \texttt{Bird} with \texttt{speak()}, subclasses \texttt{Parrot} and \texttt{Crow} that override it, and a function \texttt{make_bird_speak()} that accepts any bird object[cite: 3].

\begin{lstlisting}[language=Python]
class Bird:
    def speak(self):
        return "Chirp"

class Parrot(Bird):
    def speak(self):
        return "Squawk"

def make_bird_speak(bird_obj):
    # This polymorphically invokes .speak() without checking the object's specific class type
    print(bird_obj.speak())

# --- Execution ---
make_bird_speak(Parrot()) # Output: Squawk
\end{lstlisting}

\noindent\textbf{Detailed Explanation and Commentary:}
\textbf{This highlights duck typing}, where our consumer function relies entirely on the presence of the \texttt{.speak()} method rather than checking object classes explicitly[cite: 3]. Polymorphism lets us write flexible, modular code that handles new subclasses effortlessly[cite: 3].

\subsection*{Task 2: Polymorphic Vehicle Movement}
\textbf{Task:} Create two classes \texttt{Car} and \texttt{Bicycle} with a \texttt{move()} method, store them in a list, and loop through them calling \texttt{move()}[cite: 3].

\begin{lstlisting}[language=Python]
class Car:
    def move(self):
        return "Car drives"

class Bicycle:
    def move(self):
        return "Bicycle pedals"

# --- Execution ---
vehicles = [Car(), Bicycle()]
for v in vehicles:
    print(v.move())
\end{lstlisting}

\noindent\textbf{Detailed Explanation and Commentary:}
\textbf{This demonstrates polymorphic iteration}, allowing us to call identical method names across completely different classes without verifying their specific types beforehand[cite: 3].

\section*{Course 5: Variable Scope 🗺️✨}
Scope dictates where variables live and can be accessed within our code, functioning much like personal drawers versus shared household tables.

\subsection*{Task 1: Modifying Global State}
\textbf{Task:} Create a global variable \texttt{total = 0}, and inside a function \texttt{add_to_total(n)}, use the \texttt{global} keyword to modify it across multiple calls[cite: 3].

\begin{lstlisting}[language=Python]
# Initialize our global total tracker
total = 0

def add_to_total(n):
    # This declares that we want to modify the global 'total' variable directly
    global total
    total += n
    print("Updated global total:", total)

# --- Execution ---
add_to_total(5)
add_to_total(10)
\end{lstlisting}

\noindent\textbf{Detailed Explanation and Commentary:}
\textbf{This demonstrates how state persistence} is achieved using the \texttt{global} keyword[cite: 3]. Ordinarily, variables created inside functions are temporary; adding \texttt{global total} instructs Python to link our function directly to the module-level variable to accumulate running sums[cite: 3].

\subsection*{Task 2: Nonlocal Keyword in Nested Functions}
\textbf{Task:} Create a function \texttt{outer()} with a variable \texttt{message = "Hi"}, and inside it, create a nested function \texttt{inner()} that modifies \texttt{message} using the \texttt{nonlocal} keyword[cite: 3].

\begin{lstlisting}[language=Python]
def outer():
    message = "Hi"
    
    def inner():
        # This modifies the enclosing variable from the parent scope
        nonlocal message
        message = "Hello from inner scope!"
        
    inner()
    print("Updated outer message:", message)

# --- Execution ---
outer()
\end{lstlisting}

\noindent\textbf{Detailed Explanation and Commentary:}
\textbf{This illustrates how the \texttt{nonlocal} keyword} enables nested inner functions to mutate variables residing in their parent enclosing function without elevating them to global status[cite: 3].

\section*{Assignment \& Assessment Highlights 🚀📚}

\subsection*{Task 1: Instance Tracking with Class Variables}
\textbf{Task:} Create a \texttt{Counter} class with a class variable to track total instances created, and print the count after making several objects[cite: 3].

\begin{lstlisting}[language=Python]
class Counter:
    # Class-level variable shared across all instances
    total_instances = 0
    
    def __init__(self):
        Counter.total_instances += 1

# --- Execution ---
obj1 = Counter()
obj2 = Counter()
obj3 = Counter()
print("Total instances created:", Counter.total_instances)
\end{lstlisting}

\noindent\textbf{Detailed Explanation and Commentary:}
\textbf{This highlights class-level variables}, which maintain a shared state across all generated instances to track cumulative counts rather than resetting per object[cite: 3].

\subsection*{Task 2: Polymorphic Shape Area Aggregation}
\textbf{Task:} Create classes \texttt{Rectangle}, \texttt{Circle}, and \texttt{Triangle} with an \texttt{area()} method, store instances in a list, and calculate total area without checking types[cite: 3].

\begin{lstlisting}[language=Python]
import math

class Rectangle:
    def __init__(self, w, h):
        self.w, self.h = w, h
    def area(self):
        return self.w * self.h

class Circle:
    def __init__(self, r):
        self.r = r
    def area(self):
        return math.pi * (self.r ** 2)

# Store mixed shape instances in a single collection
shapes = [Rectangle(3, 4), Circle(2)]

# Polymorphically sum up areas without checking individual object types
total_area = sum(shape.area() for shape in shapes)
print(f"Calculated Total Area: {total_area:.2f}")
\end{lstlisting}

\noindent\textbf{Detailed Explanation and Commentary:}
\textbf{This illustrates polymorphic aggregation}, enabling uniform processing of diverse object structures[cite: 3]. Because every shape implements its own \texttt{.area()} method under a shared interface, our summation loop evaluates items seamlessly without conditional type checks[cite: 3].

\vspace{0.5cm}
\noindent What an incredible week of growth and discovery at DataraFlow! 🤗 I hope this deep dive into our exercises helps illuminate these core Python concepts for your own coding journey. Let me know which topic you would like us to explore next! ❤️✨

\end{document}